%%%%%%%%%%%%%%%%%%%%%%%%%%%%%%%%%%%%%%%%%%%%%%%%%%%%%%%%%%%%%%%%%%%
%                                                                 %
%   EPIO - Reference Manual -- LaTeX Source                       %
%                                                                 %
%   Front Material: Title page,                                   %
%                   Copyright Notice                              %
%                   Preliminary Remarks                           %
%                   Table of Contents                             %
%   EPS file      : cern15.eps, cnastit.eps                       %
%                                                                 %
%   Editor: Michel Goossens / CN-AS                               %
%   Last Mod.: 18 Nov 1993 22:45 mg                               %
%                                                                 %
%%%%%%%%%%%%%%%%%%%%%%%%%%%%%%%%%%%%%%%%%%%%%%%%%%%%%%%%%%%%%%%%%%%

%%%%%%%%%%%%%%%%%%%%%%%%%%%%%%%%%%%%%%%%%%%%%%%%%%%%%%%%%%%%%%%%%%%%
%    Tile page                                                     %
%%%%%%%%%%%%%%%%%%%%%%%%%%%%%%%%%%%%%%%%%%%%%%%%%%%%%%%%%%%%%%%%%%%%
\notHTML{\def\Ptitle#1{\special{ps: /Printstring (#1) def}
                       \epsfbox{/user/goossens/cnasall/cnastit.eps}}}
\HTML{\def\Ptitle#1{<B>#1</B>}}
 
\begin{titlepage}
\notHTML{\vspace*{-23mm}}%
\notHTML{\mbox{\epsfig{file=/usr/local/lib/tex/ps/cern15.eps,height=30mm}}}%
\HTML{<IMG SRC="../cernlogo.gif">}%
\hfill
\raise8mm\hbox{\Large\bf CERN Program Library Long Writeup I101}
\hfill\mbox{}
\HTML{<P>}
\begin{center}
\mbox{}\\[10mm]
\mbox{\Ptitle{EPIO}}\\[2cm]
\HTML{<P>\\}
{\LARGE Experimental Physics Input Output Package}\\[2cm]
\HTML{<P>\\}
{\LARGE User's Guide}\\[3cm]
\HTML{<P>\\}
{\Large Application Software and Databases Group}\\[1cm]
{\Large Computing and Networks Division}\\[2cm]
\end{center}
\HTML{<P>}
\notHTML{\vfill}%
\begin{center}\Large CERN Geneva, Switzerland\end{center}
\end{titlepage}

\Filename{H1Preface}
\HTML{<H1>Preface</H1>}

%%%%%%%%%%%%%%%%%%%%%%%%%%%%%%%%%%%%%%%%%%%%%%%%%%%%%%%%%%%%%%%%%%%%
%    Copyright  page                                               %
%%%%%%%%%%%%%%%%%%%%%%%%%%%%%%%%%%%%%%%%%%%%%%%%%%%%%%%%%%%%%%%%%%%%
\HTML{\PRE}
\thispagestyle{empty}
\framebox[\textwidth][t]{\hfill\begin{minipage}{0.96\textwidth}%
\vspace*{3mm}\begin{center}Copyright Notice\end{center}
\parskip\baselineskip
{\bf EPIO --- Experimental Physics Input Output Package}
 
CERN Program Library entry {\bf I101}
 
\copyright{} Copyright CERN, Geneva 1993
 
Copyright and any other appropriate legal protection of these
computer programs and associated documentation reserved in all
countries of the world.
 
These programs or documentation may not be reproduced by any
method without prior written consent of the Director-General
of CERN or his delegate.
 
Permission for the usage of any programs described herein is
granted apriori to those scientific institutes associated with
the CERN experimental program or with whom CERN has concluded
a scientific collaboration agreement.
 
Requests for information should be addressed to:
\vspace*{-.5\baselineskip}
\begin{center}
\tt\begin{tabular}{l}
CERN Program Library Office              \\
CERN-CN Division                         \\
CH-1211 Geneva 23                        \\
Switzerland                              \\
Tel.      +41 22 767 4951                \\
Fax.      +41 22 767 7155                \\
Bitnet:   CERNLIB@CERNVM                 \\
DECnet:   VXCERN::CERNLIB (node 22.190)  \\
Internet: CERNLIB@CERNVM.CERN.CH
\end{tabular}
\end{center}
\vspace*{2mm}
\end{minipage}\hfill}%end of minipage in framebox
\vspace{6mm}
\HTML{<P>}
 
{\bf Trademark notice: All trademarks appearing in this guide are acknowledged as such.}
\vfill
\HTML{<P>}
\begin{tabular}{l@{\quad}l@{\quad}>{\tt}l}
{\em Contact Person\/}:        & Ian McLaren /CN     & (mclareni\atsign cernvm.cern.ch)\\[1mm]
{\em Technical Realization\/}: & Michel Goossens /CN & (goossens\atsign cernvm.cern.ch)\\[5mm]
{\em Edition -- December 1993}
\end{tabular}
\HTML{\ePRE}%
\newpage
 
%%%%%%%%%%%%%%%%%%%%%%%%%%%%%%%%%%%%%%%%%%%%%%%%%%%%%%%%%%%%%%%%%%%%
%    Introductory material                                         %
%%%%%%%%%%%%%%%%%%%%%%%%%%%%%%%%%%%%%%%%%%%%%%%%%%%%%%%%%%%%%%%%%%%%
\pagenumbering{roman}
\setcounter{page}{1}

\Filename{H2EPIO-History}
\section*{History}

Hans Grote/SL was one of the original two authors of the EPIO package. 

The EPIO offline package is intended to be used for all data (at
least on tape) in an experiment, in such a way that from the raw data
tape to the DST, the tape (or file) format is identical. This makes the
transport of data between computers trivial, and simplifies the task of
passing the files or tapes at different stages of the production chain
through any other part of the production chain.
 
The package has been developed at CERN as a part of \Lit{packlib}.
\index{packlib@\Lit{packlib}}
 
Version 1.66 of EPIO introduced a standard Fortran 77 version which is
the base for the Unix versions.

We would like to encourage all users to contact us in case:
 
\begin{itemize}
\item  they detect mistakes in the manual or the package,
\item  they would like to see something to be added to the
       manual or described in more detail,
\item  they would like to see features added to the package.
\end{itemize}
 
Remember that your collaboration in this is vital, and that the poor
man in (Saclay, Rome, Vienna, Hamburg, etc.) may desperately be trying to
find and fix a bug that you have just corrected, but told nobody (and
vice versa, of course).

\Filename{H2Preliminary-remarks}
\section*{Preliminary remarks}
 
This manual serves at the same time as a {\bf Reference manual}
and as a {\bf User Guide} for the EPIO package.
 
In this manual
examples are in {\tt monotype face} and strings to be input by the user 
are {\tt\underline{underlined}}.
In the index the page where a routine is defined is in {\bf bold},
page numbers where a routine is referenced are in normal type.

In the description of the routines a \Lit{*} following
the name of a parameter indicates that this is an {\bf output} parameter.
If another \Lit{*} precedes a parameter in the calling sequence, the
parameter in question is both an {\bf input} and {\bf output} parameter.

This document has been produced using \LaTeX~\cite{bib-LATEX}
with the \Lit{cernman} style option, developed at CERN. 
A compressed PostScript file \Lit{epio.ps.Z}, 
containing a complete printable version
of this manual, can be obtained from any CERN machine
by anonymous ftp as follows
(commands to be typed by the user are underlined):

\vspace*{3mm} 
\begin{XMP}
    \underline{ftp asis01.cern.ch}
    Trying 128.141.201.136...
    Connected to asis01.cern.ch.
    220 asis01 FTP server (Version 6.10 ...) ready.
    Name (asis01:username): \underline{anonymous}
    Password: \underline{your\_{}mailaddress}
    230 Guest login ok, access restrictions apply.
    ftp> \underline{cd cernlib/doc/ps.dir}
    ftp> \underline{binary}
    ftp> \underline{get epio.ps.Z}
    ftp> \underline{quit}
\end{XMP}

%%%%%%%%%%%%%%%%%%%%%%%%%%%%%%%%%%%%%%%%%%%%%%%%%%%%%%%%%%%%%%%%%%%%
%    Tables of contents ...                                        %
%%%%%%%%%%%%%%%%%%%%%%%%%%%%%%%%%%%%%%%%%%%%%%%%%%%%%%%%%%%%%%%%%%%%
\newpage
\tableofcontents
%\newpage
%\listoffigures
\listoftables

